\documentclass[letterpaper]{article}

\usepackage[letterpaper,margin=1in]{geometry}

% Compile from the repository root so that this path resolves to
% the pinned WSP submodule.
\input{wsp/testing/test-case-library.tex}

\begin{document}

\def\TCID{%
  TC-0001
}

\def\TCTitle{%
  Member table of the common controls
}

\def\TCPurpose{%
  Verify that the generated member table lists the nine MSComctlLib classes in
  order and that every class has named members of a known kind and events that
  are methods.
}

\def\TCPriority{%
  \TCPriorityHigh
}

\def\TCRequirementRef{%
  REQ-0001, REQ-0002, REQ-0003, REQ-0004, REQ-0005, REQ-0006, REQ-0007,
  REQ-0008, REQ-0009
}

\def\TCDesignRef{%
  Not applicable; the rules are in the named header.
}

\def\TCFeatureRef{%
  \texttt{src/mscomctl\_tables.cpp},
  \texttt{include/wfc/mscomctl/mscomctl.hpp}
  (\texttt{wfc::ui::mscomctl\_classes()}) and
  \texttt{tests/mscomctl\_tables\_test.cpp} and the runner
  \texttt{tests/tc-0001-tables.ps1}
}

\def\TCTestTechnique{%
  Equivalence partitioning and boundary value analysis of the rule functions
  (ISO/IEC/IEEE 29119-4 specification-based techniques).
}

\def\TCPreconditions{%
  The library is built standalone (\texttt{cmake -S . -B out/build}, then
  \texttt{cmake --build out/build --config Debug}) with
  \texttt{WFC\_BUILD\_TESTS} on; CTest runs the test as
  \texttt{TC-0001-tables}
  (or \texttt{wfc-mscomctl-TC-0001-tables} when WFC adds
  the library).
}

\def\TCEnvironment{%
  Any target the library builds for (Windows x86, x64 and ARM64 with MSVC,
  Linux through Winelib); local runs use the x64 Debug configuration of the
  standalone build. The test needs no window, registry or network.
}

\def\TCAssumptions{%
  The rule header under test is the one built into the library; the test is
  compiled from the same source revision.
}

\def\TCInputData{%
  The table itself; no other input.
}

\def\TCInitialState{%
  A fresh process; the test keeps no state between runs and holds no files or
  handles.
}

\def\TCProcedure{%
  \begin{enumerate}
    \item Run \texttt{mscomctl\_tables\_test} (CTest \texttt{TC-0001-tables}).
    \item Read the class list through \texttt{wfc::ui::mscomctl\_classes()}
      and, for each class, every member row and event row.
  \end{enumerate}
}

\def\TCExpectedResult{%
  \begin{enumerate}
    \item The table has exactly nine classes: TabStrip, Toolbar, StatusBar,
      ProgressBar, TreeView, ListView, ImageList, Slider, ImageCombo, in that
      order.
    \item Every class has at least one member; every member has a non-empty
      name and a kind of \texttt{g}, \texttt{l}, \texttt{s} or \texttt{m};
      every event row has the kind \texttt{m}.
    \item The program prints \texttt{mscomctl: 9 classes} and exits with code
      0.
  \end{enumerate}
}

\def\TCPostConditions{%
  The process ends; nothing is written outside CTest's own log.
}

\TCTemplate

\end{document}
